%------------------------------------------------------------------------------------
\newpage
%\chapter{ТЕОРЕТИЧЕСКАЯ ЧАСТЬ}
\begin{center}
  \textbf{\large 1. ТЕОРЕТИЧЕСКАЯ ЧАСТЬ}
\end{center}
\refstepcounter{chapter}
\addcontentsline{toc}{chapter}{1. ТЕОРЕТИЧЕСКАЯ ЧАСТЬ}

\section{ OFDM модуляция}

Данный вид модуляции подразумевает передачу суммы нескольких гармоник
— поднесущих, частоты которых выбираются из условия ортогональности [\ref{Dahlman2011}]
. Математическое представление OFDM сигнала приведено в следующем выражении (\ref{eq:ofdm_base}):

%------------------------------------------------------------------------------------
% Equation
%------------------------------------------------------------------------------------
\begin{equation}
    S(t) = \sum_{n = 0}^{N - 1} X[n] e^{j 2 \Pi f[n] t},
    \label{eq:ofdm_base}
\end{equation}
%------------------------------------------------------------------------------------

\noindent где $S(t)$ - OFDM символ, $N$ - количество поднесущих, $X[n]$ - отсчет полезного сигнала,
$f[n]$ - частота $n$-ой поднесущей.

Частоты поднесущих получаются из выражения (\ref{eq:subc_eq}):

%------------------------------------------------------------------------------------
% Equation
%------------------------------------------------------------------------------------
\begin{equation}
    f[n] = n \Delta f,
    \label{eq:subc_eq}
\end{equation}
%------------------------------------------------------------------------------------

\noindent где $\Delta f$ - расстояние между поднесущими.

%------------------------------------------------------------------------------------
% Figure
%------------------------------------------------------------------------------------
\begin{figure}[h]
    \centering
    \includegraphics[width=0.8\textwidth]{images/1_ofdm_fd.pdf}
    \caption{OFDM сигнал в частотной области}
    \label{fig:ofdm_fd}
\end{figure}
\FloatBarrier
%------------------------------------------------------------------------------------

На рис.~\ref{fig:ofdm_fd} показан OFDM сигнал в частотной области, где $k$ - индекс поднесущей.
Из графика видно, что отдельные поднесущие не перекрываются из-за ортогональности.

При последовательной передаче высокоскоростного потока данных, как
правило, возникает проблема, связанная c тем, что период передачи символов
намного меньше, чем время задержки в канале [\ref{Sesia2009}]
.
Это создает помехи в виде
символьной интерференции, в ходе чего происходит наложение символов, которые можно устранить c помощью эквализации.
В OFDM модуляции высокоскоростной поток символов сначала последовательно-параллельным
образом преобразуется для модуляции на $N$ параллельных
поднесущих (рис.~\ref{fig:ofdm_tx}). Это увеличивает период передачи символов на каждой
поднесущей так, что он становится значительно больше, чем временные
задержки в канале. После модуляции, сигнал поступает
на блок ОБПФ и на параллельно-последовательный преобразователь, на котором вводится циклический
префикс - введение избыточности в начале OFDM символа для борьбы с межсимвольной интерференцией.
Приемник устроен аналогично (рис.~\ref{fig:ofdm_rx}).

%------------------------------------------------------------------------------------
% Figure
%------------------------------------------------------------------------------------
\begin{figure}[h]
    \centering
    \includegraphics[width=0.8\textwidth]{images/1_ofdm_tx.png}
    \caption{Схема передатчика OFDM}
    \label{fig:ofdm_tx}
\end{figure}
\FloatBarrier
%------------------------------------------------------------------------------------

%------------------------------------------------------------------------------------
% Figure
%------------------------------------------------------------------------------------
\begin{figure}[h]
    \centering
    \includegraphics[width=0.8\textwidth]{images/1_ofdm_rx.png}
    \caption{Схема приемника OFDM}
    \label{fig:ofdm_rx}
\end{figure}
\FloatBarrier
%------------------------------------------------------------------------------------


%------------------------------------------------------------------------------------
\section{Эффект Доплера}

Эффект Доплера возникает при наличии радиальной
скорости приемника или передатчика относительно друг друга (\ref{eq:dopp}). 

%------------------------------------------------------------------------------------
% Equation
%------------------------------------------------------------------------------------
\begin{equation}
    f_d = \frac{v f}{c} \cos \psi ,
    \label{eq:dopp}
\end{equation}
%------------------------------------------------------------------------------------

\noindent где $f_d$ - доплеровское смещение частоты, $v$ - модуль относительной скорости между
премником и передатчиком, $f$ - частота сигнала, $\psi$ - угол между вектором скорости и
направлением от передатчика к приемнику  (рис.~\ref{fig:doppler}).

%------------------------------------------------------------------------------------
% Figure
%------------------------------------------------------------------------------------
\begin{figure}[h]
    \centering
    \includegraphics[width=0.5\textwidth]{images/1_doppler.png}
    \caption{Схема для расчета доплеровского смещения}
    \label{fig:doppler}
\end{figure}
\FloatBarrier
%------------------------------------------------------------------------------------

\newpage
В канале связи при наличии отражающих и рассеивающих объектов возникает многолучевое
распространение радиосигнала, которое описывается
детерминироваными и стохастическими моделями каналов.

Если считать сигнал узкополосным, то эффект Доплера приводит только к смещению спектра. В моделях
канала с многолучевым распространением (напр. модель Джейкса), вводится такое допущение, т.е. 
смещение действует на весь спектр в целом. Вследствии многолучевости,
отраженные сигналы приходят на приемник с разных углов и с разным смещением, тем самым достигается
эффект масштабирования спектра. Такой подход позволяет достаточно просто учитывать и многолучевость,
и эффект Допплера.

Если такое допущение не принимать, т.е. считать сигнал широкополосным, то на каждую поднесущую
сигнала будет действовать разный сдвиг частоты, что приводит к интерференции между поднесущими (ICI)
без учета многолучевости. Это можно представить следующим образом (\ref{eq:ofdm_dopp_full}):

%------------------------------------------------------------------------------------
% Equation
%------------------------------------------------------------------------------------
\begin{equation}
    S(t) = \sum_{n = 0}^{N - 1} X[n] e^{j 2 \Pi (f[n] + f_c) (1 + \varepsilon ) t},
    \label{eq:ofdm_dopp_full}
\end{equation}
%------------------------------------------------------------------------------------

\noindent где $f_c$ - несущая частота сигнала, $\varepsilon = \frac{v}{c} \cos \psi$ - коэффициент доплеровского смещения.

Обычно, в спутниковых системах связи доплеровское смещение заранее известно, так как 
известны координаты и скорость спутника, поэтому можно предавать сигнал с учетом
смещения $S^`(t) = S(t) e^{- 2 j \Pi f_c \varepsilon t}$, тем самым избавляясь от этого эффекта.
В сетях наземной мобильной связи, где отсутсвует информация о координатах и скоростях приемника/передатчика,
можно по специальным сигналам для оценки канала оценить смещение и отфильтровать его
на эквалайзере. Наконец, смещение несущей частоты может исправить система автоподстойки частоты (АПЧ).
Исходя из этого, в дальнейших рассуждениях несущая частота сигнала не будет учитываться,
так как ее вклад в искажение сигнала зачастую компенсируется.

В данной работе влияние эффекта Доплера сводится к следующему выражению (\ref{eq:ofdm_dopp_real}):
%------------------------------------------------------------------------------------
% Equation
%------------------------------------------------------------------------------------
\begin{equation}
    S(t) = \sum_{n = 0}^{N-1} X[n] e^{j 2 \Pi f[n] (1 + \varepsilon ) t},
    \label{eq:ofdm_dopp_real}
\end{equation}
%------------------------------------------------------------------------------------
\noindent где $f_n = n \Delta f$, $\Delta f = \frac{1}{T}$, $T$ - длительность OFDM символа.

%------------------------------------------------------------------------------------
\section{Оценка ICI и SINR}


%------------------------------------------------------------------------------------
% Equation
%------------------------------------------------------------------------------------
%\begin{equation}
%    S(t) = \sum_{n = 0}^{N - 1} X_n e^{j 2 \Pi f_n (1 + \varepsilon ) t},
%    \label{eq:ofdm_1}
%\end{equation}
%------------------------------------------------------------------------------------

%где $f_n = n \Delta f$, $\Delta f = \frac{1}{T}$, $T$ - длительность OFDM символа.

Для нахождения ICI [\ref{Li2018}] запишем выражение для сигнала на приемнике:
%------------------------------------------------------------------------------------
% Equation
%------------------------------------------------------------------------------------
\begin{equation}
    Y[k] = \frac{1}{T} \int_{0}^{T} S(t) e^{-j 2 \Pi f[k] t} \,dt,
    \label{eq:ofdm_2}
\end{equation}
%------------------------------------------------------------------------------------

\noindent где $f[k] = k \Delta f$ - частота поднесущей для $k$-го отсчета.

\noindent Подставим (\ref{eq:ofdm_dopp_real}) в (\ref{eq:ofdm_3}):

%------------------------------------------------------------------------------------
% Equation
%------------------------------------------------------------------------------------
\begin{equation}
    Y[k] = \frac{1}{T} \sum_{n = 0}^{N - 1} X[n] \int_{0}^{T}  e^{j 2 \Pi [f[n] (1 + \varepsilon ) - f[k]] t} \,dt,
    \label{eq:ofdm_3}
\end{equation}
%------------------------------------------------------------------------------------

\noindent Рассмотрим следующий интеграл (\ref{eq:integral_1}):
%------------------------------------------------------------------------------------
% Equation
%------------------------------------------------------------------------------------
\begin{equation}
    \label{eq:integral_1}
    \begin{split}
        I[n, k] &= \frac{1}{T} \int_{0}^{T}  e^{j 2 \Pi \Delta f [n (1 + \varepsilon ) - k] t} \,dt \\
        &= \frac{1}{T} \left[ \frac{e^{j 2 \Pi \Delta f [n (1 + \varepsilon ) - k] t}}{j 2 \Pi \Delta f [n (1 + \varepsilon ) - k]} \right]_{0}^{T} \\
        &= \frac{e^{j 2 \Pi [n (1 + \varepsilon ) - k]} - 1}{j 2 \Pi [n (1 + \varepsilon ) - k]}
    \end{split}
\end{equation}

\noindent Рассмотрим выражение вида (\ref{eq:integral_2}):

\begin{equation}
    \label{eq:integral_2}
    \begin{split}
        exp(2 j \Pi \beta ) - 1 &= exp(j \Pi \beta ) exp(j \Pi \beta ) - exp(j \Pi \beta ) exp(-j \Pi \beta ) \\
        &= exp(j \Pi \beta) [exp(j \Pi \beta) - exp(- j \Pi \beta)] = 2 j exp(j \Pi \beta) \sin(\Pi \beta)
    \end{split}
\end{equation}

\noindent Подставим в (\ref{eq:integral_1}):
%------------------------------------------------------------------------------------
% Equation
%------------------------------------------------------------------------------------
\begin{equation}
    I[n, k] = exp(j \Pi [(1 + \varepsilon )n - k]) sinc(\Pi [ (1 + \varepsilon) n - k])
    \label{eq:matrix}
\end{equation}
%------------------------------------------------------------------------------------

\noindent Таким образом, сигнал на приемнике будет иметь следующий вид (\ref{eq:rx_signal}):
%------------------------------------------------------------------------------------
% Equation
%------------------------------------------------------------------------------------
\begin{equation}
    Y[k] = \sum_{n = 0}^{N - 1} X[n] exp(j \Pi [(1 + \varepsilon )n - k]) sinc(\Pi [ (1 + \varepsilon) n - k]),
    \label{eq:rx_signal}
\end{equation}
%------------------------------------------------------------------------------------

\newpage
Выражение (\ref{eq:rx_signal}) можно представить в следующем виде (\ref{eq:rx_signal_matrix}), с
учетом влияния аддитивного белого гауссовского шума:
%------------------------------------------------------------------------------------
% Equation
%------------------------------------------------------------------------------------
\begin{equation}
    Y[k] = X[k] I[k, k] + \sum_{n = 1, n \neq k}^{N} X[n] I[n, k] + n = X[k] I[k, k] + \alpha ,
    \label{eq:rx_signal_matrix}
\end{equation}
%------------------------------------------------------------------------------------
\noindent где $X[k] I[k, k]$ - полезный сигнал c поднесущей $k$, $\alpha = \sum_{n = 1, n \neq k}^{N} X[n] I[n, k] + n$
- сигнал полученный от других поднесущих в результате ICI, $n$ - белый гауссов шум.

При отсутсвии эффекта Доплера ($\varepsilon = 0$)
элементы (\ref{eq:matrix}) составляют единичную матрицу: $I[k, k] = 1$,
$I[n, k] = 0, n \neq k$, т.к.
$sinc(\Pi m) = 0$, $\forall m \in \mathbb{Z}, m \neq 0$. Таким образом реализуется передача данных в
канале без ICI под воздействием белого аддитивного гауссовского шума.

Для оценки сигналов при наличии интерференции вводят отношение полезного сигнала к интерференции и шуму
(SINR [\ref{Goldsmith2005}], Signal to Interference plus Noise Ratio):
%------------------------------------------------------------------------------------
% Equation
%------------------------------------------------------------------------------------
\begin{equation}
    SINR = \frac{P_{signal}}{P_{ICI} + P_{noise}},
    \label{eq:sinr}
\end{equation}
%------------------------------------------------------------------------------------
где $P_{signal}$ - мощность полезного сигнала, $P_{ICI}$ - мощность интерференции, $P_{noise}$
- мощность шума.

Учитывая (\ref{eq:rx_signal_matrix}), SINR будет различным для каждой поднесущей:
%------------------------------------------------------------------------------------
% Equation
%------------------------------------------------------------------------------------
\begin{equation}
    SINR[k] = \frac{ \mathbb{E} \left [|X[k] I [k, k]|^2 \right ] }
                   { \mathbb{E} \left [|\sum_{n = 1, n \neq k}^{N} X[n] I[n, k]|^2 \right ] + \sigma_n^2 },
    \label{eq:sinr_k}
\end{equation}
%------------------------------------------------------------------------------------
где $\sigma_n^2$ - дисперсия шума.

Для удобства введем значение $SINR[k]$ в дБ (\ref{eq:sinr_k_db}):
%------------------------------------------------------------------------------------
% Equation
%------------------------------------------------------------------------------------
\begin{equation}
    SINR_{dB}[k] = 10 \log_{10} \frac{ \mathbb{E} \left [|X[k] I [k, k]|^2 \right ] }
                   { \mathbb{E} \left [|\sum_{n = 1, n \neq k}^{N} X[n] I[n, k]|^2 \right ] + \sigma_n^2 }
    \label{eq:sinr_k_db}
\end{equation}
%------------------------------------------------------------------------------------

На рис.~\ref{fig:sinr_theor} представлены кривые $SINR_{dB}$ (\ref{eq:sinr_k_db}) от индекса поднесущей
$k$ при различных заначениях модуля скорости относительного движения приемника/передатчика.

%------------------------------------------------------------------------------------
% Figure
%------------------------------------------------------------------------------------
\begin{figure}[h]
    \centering
    \includegraphics[width=1\textwidth]{images/1_theor_sinr.pdf}
    \caption{Теоретические кривые для зависимости $SINR_{dB}$ от $N$}
    \label{fig:sinr_theor}
\end{figure}
\FloatBarrier
%------------------------------------------------------------------------------------

%------------------------------------------------------------------------------------
\section{5G NR NTN}


%------------------------------------------------------------------------------------
\section{Орбиты}


%------------------------------------------------------------------------------------
\section{Модели канала}

% Здесь тоже можно оставить пару ссылок на работы

% Нужно сказать, что 3gpp порекомендовали CDL и TDL для работы в частотном диапазоне [a, b]
% В такой то статье

Для моделирования каналов связи существует множество моделей канала. Для исследований в области 
систем связи 3gpp выпустили документ с рекомендованными к использованию моделями канала в частотном диапазоне
от 0.5 МГц до 100 МГц [Ссылка на документ].
%Их можно разделить на стохастические и детерминированные.

% Здесь нужно дать подробную классификацию моделей канала

В исследованиях спутниковой связи используются преимущественно две из них 
- TDL (Tapped Delay Line) и СDL (Cluster Delay Line). 
Обе эти модели подразумевают распростанение сигнала по нескольким пространственным путям. % тут можно вставить картинку

TDL является 2D моделью и не учитывает обьемное распространение сигнала, диаграммы направленностей
антенн, в то время как
 СDL модель канала используется для учета распрастранения сигнала в обьемном канале (3d),
 учета диаграмм  направленностей, поляризации и тд.

Далее будет дана краткая характеристика каналов связи, их основные параметры

% Не нужно полностью описывать устройство CDL модели, которую мы не будем использовать


\subsection{TDL}
% Ссылка на работу 
% On the influence of Delay Spread in TDL and CDL Channel Models for Downlink 5G MIMO Systems

Модель TDL подразумевает распространение сигнала по нескольким путям (taps).

% Картинку с TDL моделью 

Испульсная характеристика канала описывается выражением (\ref{eq:tdl_impulse_response}):

%------------------------------------------------------------------------------------
% Equation
%------------------------------------------------------------------------------------
\begin{equation}
    H(t, \tau) = \sum_{i = 0}^{n} a_i(t) \delta (\tau - \tau_i),
    \label{eq:tdl_impulse_response}
\end{equation}
%------------------------------------------------------------------------------------
\noindent где $a_i(t)$ - амплитуда $\tau_i$ задержки на $i$ом пути распространения сигнала.

% Для NTN существуют другие профайлы, которые можно найти в следующей статье
В спецификации существуеют 5 различных наборов параметров TDL для различных сценариев:
\begin{itemize}
    \item Для ситуации с отсутствием прямой видимости (NLOS, Non-Line-of-Side)
        \begin{itemize}
            \item TDL_A
            \item TDL_B
            \item TDL_C
        \end{itemize}

    \item Для ситуации с наличием прямой видимости (LOS, Line-of-Side)
        \begin{itemize}
            \item TDL_D
            \item TDL_E
        \end{itemize}
\end{itemize}

% Ссылка на Doppler effect mitigation in LEO-based 5G Non-Terrestrial Networks
Для моделирования будет использоваться модель канала из работы (\ref{}), которая основана на модели TDL.

Импульсная характеристика меняющимейся во времени частотно-селективной модели представлена как (\ref{}):

% Данное выражение нужно исправить и понять, в чем тут соль

%------------------------------------------------------------------------------------
% Equation
%------------------------------------------------------------------------------------
\begin{equation}
    р(t, \tau)^{(i, j)} = \sqrt{P_p(t)} \sum_{q = 0}^{Q_p(t) - 1} a_{p,q}^{(i,j)}(t)
                          \exp^{j 2 Pi f_{D, p, q}(t) t} \delta(\tau - \tau_{p, q}(t)),
    \label{eq:real_channel_impulse_response}
\end{equation}
%------------------------------------------------------------------------------------
\noindent где $i$ прин. ${0, 1, ..., N_{RX} - 1}$, $j$ прин. ${0, 1, ..., N_{TX} - 1}$,
$P_p = 10^{-(P_d + P_S)/10}$, $P_d = 20log_{10} (4 \Pi d / \lambda)$ отвечает за 
потери сигнала в пространстве при прохождении растояния $d$ между спутником и $p-th$ UE,
$P_S ~ N(0, \gamma^2_S)$ относится к затемнению (shadowing) с \gamma^2_S дисперсией.
$q$, $Q_p$, $a_{p,q}$, $\theta_{p,q}$ обозначают индекс пути, количество путей, амплитуды этих путей и 
задержки для $p-th$ UE и $q$го пути. 
% Добавить сверху формулы для доплера, которые будут использоваться в данном случае


%\subsection{SDL}
% Здесь можно оставить ссылку на тот обзор на русском языке, который ты нашел

